% Requires: \usepackage{xcolor,booktabs,tabularx,array}
\definecolor{normalgray}{HTML}{4F5965}
\definecolor{wideblue}{HTML}{2457A6}
\definecolor{netgreen}{HTML}{178A62}
\definecolor{deeporange}{HTML}{D66B20}
\begin{table*}[t]
\raggedright
\scriptsize
\renewcommand{\arraystretch}{1.05}
\begin{tabularx}{\textwidth}{*{4}{>{\raggedright\arraybackslash}X}}
\toprule
\textcolor{normalgray}{\textbf{\large Normal play}} & \textcolor{wideblue}{\textbf{\large Wide front-court switching with extended rallies}} & \textcolor{netgreen}{\textbf{\large Aggressive short-court net play}} & \textcolor{deeporange}{\textbf{\large Deep back-court defensive reset}} \\
\addlinespace[2pt]
No tactical instruction; follow the learned normal policy. & Target the wide front corners and switch depth frequently to force long, wide exchanges. & Attack aggressively with short, sharp shots to the front court to force errors and minimize movement. & Focus on deep back-court placements to keep the opponent at bay and avoid wide or short shots. \\
\addlinespace[2pt]
\textit{\textbf{Starting situation.} Source match 3, set 2, rally 13. Kento MOMOTA serves against CHOU Tien Chen at 6--6. Their normalized court positions are (-0.01, 0.50) and (-0.22, 0.17), respectively. Every case starts from this identical situation with an empty rally history.} & \textit{\textbf{Starting situation.} Source match 3, set 2, rally 13. Kento MOMOTA serves against CHOU Tien Chen at 6--6. Their normalized court positions are (-0.01, 0.50) and (-0.22, 0.17), respectively. Every case starts from this identical situation with an empty rally history.} & \textit{\textbf{Starting situation.} Source match 3, set 2, rally 13. Kento MOMOTA serves against CHOU Tien Chen at 6--6. Their normalized court positions are (-0.01, 0.50) and (-0.22, 0.17), respectively. Every case starts from this identical situation with an empty rally history.} & \textit{\textbf{Starting situation.} Source match 3, set 2, rally 13. Kento MOMOTA serves against CHOU Tien Chen at 6--6. Their normalized court positions are (-0.01, 0.50) and (-0.22, 0.17), respectively. Every case starts from this identical situation with an empty rally history.} \\
\midrule
1. Player: Serve short to Front Center & 1. Player: Serve short to Front Center & 1. Player: Serve short to Front Center & 1. Player: Serve short to Front Center \\
2. Opponent: Push Shot to Rear Right Side & 2. Opponent: Push Shot to Rear Center & 2. Opponent: Push Shot to Rear Right Side & 2. Opponent: Push Shot to Rear Right Side \\
3. Player: Smash to Mid Center & 3. Player: \textcolor{wideblue}{\textbf{Drop to Front Right Side}} & 3. Player: Clear to Rear Right Side & 3. Player: \textcolor{deeporange}{\textbf{Clear to Rear Left Side}} \\
4. Opponent: Smash Defence to Mid Center & 4. Opponent: Net Shot to Front Left Side & 4. Opponent: Smash to Mid Left Side & 4. Opponent: Clear to Rear Right Side \\
5. Player: Smash Defence to Mid Center & 5. Player: \textcolor{wideblue}{\textbf{Net Shot to Front Right Side}} & 5. Player: \textcolor{netgreen}{\textbf{Smash Defence to Front Left Side}} & 5. Player: Drop to Front Right Side \\
6. Opponent: Push Shot to Mid Center & 6. Opponent: Net Shot to Front Left Side & 6. Opponent: Net Shot to Front Right Side & 6. Opponent: Net Shot to Front Left Side \\
7. Player: Smash Defence to Front Center & 7. Player: \textcolor{wideblue}{\textbf{Lob to Rear Center}} & 7. Player: \textcolor{netgreen}{\textbf{Net Shot to Front Left Side}} & 7. Player: \textcolor{deeporange}{\textbf{Lob to Rear Left Side}} \\
8. Opponent: Net Shot to Front Left Side & 8. Opponent: Smash to Rear Right Side & 8. Opponent: Lob to Deep Right Corner & 8. Opponent: Clear to Rear Right Side \\
9. Player: Push Shot to Deep Left Corner & 9. Player: \textcolor{wideblue}{\textbf{Smash Defence to Front Center}} &  & 9. Player: Smash to Mid Right Side \\
 & 10. Opponent: Net Shot to Front Center &  & 10. Opponent: Smash Defence to Front Left Side \\
 & 11. Player: Push Shot to Mid Left Side &  & 11. Player: \textcolor{deeporange}{\textbf{Push Shot to Rear Right Side}} \\
 & 12. Opponent: Smash to Mid Left Side &  & 12. Opponent: Clear to Rear Left Side Line \\
 & 13. Player: \textcolor{wideblue}{\textbf{Smash Defence to Front Right Side}} &  &  \\
 & 14. Opponent: Push Shot to Rear Center &  &  \\
 & 15. Player: \textcolor{wideblue}{\textbf{Clear to Rear Right Side}} &  &  \\
 & 16. Opponent: Smash to Mid Center &  &  \\
 & 17. Player: \textcolor{wideblue}{\textbf{Smash Defence to Front Left Side}} &  &  \\
 & 18. Opponent: Push Shot to Rear Right Side &  &  \\
 & 19. Player: Smash to Mid Left Side &  &  \\
 & 20. Opponent: Smash Defence to Mid Right Side &  &  \\
 & 21. Player: Missed shot &  &  \\
\midrule
\textbf{Instruction alignment.} \textcolor{normalgray}{Reference trajectory without a tactical instruction.} & \textbf{Instruction alignment.} \textcolor{wideblue}{Strong alignment: four wide front-court landings and seven depth switches.} & \textbf{Instruction alignment.} \textcolor{netgreen}{Partial alignment: the fifth and seventh shots stay in the front court, and the seventh is an attacking Net Shot.} & \textbf{Instruction alignment.} \textcolor{deeporange}{Partial alignment: the third-shot Clear and seventh-shot Lob reset play to the back court; the eleventh shot also stays deep.} \\
\addlinespace[2pt]
\textbf{Result.} CHOU Tien Chen wins after Kento MOMOTA's Push Shot lands out of bounds. & \textbf{Result.} CHOU Tien Chen wins after Kento MOMOTA misses the shuttle. & \textbf{Result.} Kento MOMOTA wins after CHOU Tien Chen's Lob lands out of bounds. & \textbf{Result.} Kento MOMOTA wins after CHOU Tien Chen's Clear lands out of bounds. \\
\bottomrule
\end{tabularx}
\caption{Four StateRAG trajectories from one identical starting situation. The selected cases provide interpretable instruction-following behavior and varied terminal outcomes, including a missed shot. Colored action phrases mark Player shots or short sequences that align with the instruction in the same column; uncolored actions provide context. Player denotes Kento MOMOTA and Opponent denotes CHOU Tien Chen.}
\label{tab:kento-explanatory-case-study}
\end{table*}
